\begin{lemma}\label{lem:decic60-quintic-descent}
Let $F$ be a quintic number field, and suppose that $E/F$ has a point
of order $30$. Put $L=F(E[2])$, and suppose that $[L:F]=2$.
No curve over $L$ obtained from $E$ by a $2$-isogeny has a point
of order $60$.
\end{lemma}
\begin{proof}
The classification of quintic points on $X_1(30)$ gives the two
curves $E_{30},E_{30}'$ over
$F=\mathbb Q(\alpha)$, where
$\alpha^5+\alpha^4-3\alpha^3+3\alpha+1=0$.
We use the printed Tate models in the quintic classification.
They are $2$-isogenous over $F$.
For each model, its $2$-division polynomial has one linear and
one irreducible quadratic factor $h$ over $F$.
The exact calculations in \path{quintic_2tower.m} give
\[
 -\operatorname{disc}(h)\notin F^{\times2}.
\]
Since $[F:\mathbb Q]$ is odd, $i\notin F$; hence $i\notin L$.
The same calculation shows that the polynomial for halving the
chosen point of order $30$ is irreducible of degree $4$ over $F$,
for both curves.

Both full $2$-division fields have a degree-one prime above $53$
at which the curve has good reduction with exactly $60$ points.
The explicit quintic and quadratic residue roots are recorded in
the pinned full-$2$ certificates. Thus torsion on every
$L$-rational $2$-isogenous curve has order dividing $60$.
In particular, $E(L)_{\mathrm{tors}}=C_2\times C_{30}$.
The quotient by the $F$-rational point of order $2$ is the other
quintic model. It has no point of order $4$ over $L$: otherwise its
$2$-primary torsion is cyclic of order $4$, and halving a generator
of its $F$-rational subgroup of order $30$ would be possible over
$L$, contrary to the irreducible quartic.

Complete the square and write the roots of the $2$-division cubic
as $r_0,r_1,r_2$, with $r_0\in F$ and $r_1,r_2$ exchanged by the
nontrivial element $\sigma$ of $\operatorname{Gal}(L/F)$.
A point $(r_0,0)$ is divisible by $2$ precisely when both
$r_0-r_1$ and $r_0-r_2$ are squares. Since they are conjugate and
$E(L)$ has no point of order $4$, both are nonsquares.
The absence of order $4$ on the quotient by $(r_0,0)$ similarly
shows that $r_1-r_0$ and $r_2-r_0$ are nonsquares.
Here we use the usual $2$-isogeny formula: on the quotient by
$(r_j,0)$, the nonzero dual-kernel point is divisible by $2$ if and
only if at least one of the two differences $r_k-r_j$, $k\ne j$,
is a square.

Suppose that the quotient by $(r_1,0)$ has a point of order $4$.
Its $2$-primary torsion has order at most $4$, so that point halves
the dual-kernel point. The formula just recalled implies that
$r_2-r_1$ is a square, since $r_0-r_1$ is not.
Applying $\sigma$ shows that its negative is also a square.
This gives $i\in L$, a contradiction. The quotient by $(r_2,0)$
is conjugate and is excluded as well. None of the three quotients
has a point of order $4$, and hence none has a point of order $60$.
\end{proof}
